\section{The weave}\label{sec:weave}

\paragraph{The question.}
The grammar allows many descriptions, not one. Section~\ref{sec:grammar} selected one of them as a root, under four
inputs. The honest form of the question ``what is forced'' is what holds for all the descriptions at once, with none
picked by hand. This section names the descriptions, names their joint action, and finds the first structure that the
joint action forces: the records' own parities, which move as three objects and which give a triplet. The program's
arcs are cited by number~\cite{OriginAxiomRepo}.

\paragraph{Threads.}
A \emph{thread} is a closed path of the moves through the shared records. Concretely it is a primitive hyperbolic
conjugacy class $\pm w$ in $\SL(2,\Z)$, where $w$ is a cyclic word in $L$ and $R$ containing both letters; such classes
exhaust the hyperbolic classes (Section~\ref{sec:grammar}). A thread is also a closed geodesic on the modular surface
$\Hb/\PSL(2,\Z)$, and the word is its cutting sequence~\cite{Series1985,Moeckel1982}. Its once-punctured-torus bundle is
a cusped hyperbolic $3$-manifold. The root $LR$ of Theorem~\ref{thm:seed} is one thread, of trace $3$.

\begin{definition}[the weave]\label{def:weave}
The \emph{weave} is the joint action of every allowed move on the records, together with everything that joint action
forces. A \emph{weave statement} takes every thread under a rule stated in advance, with none chosen, is defined by the
joint action, and is a claim about all threads at once. A \emph{thread statement} is anything read on one thread or on
its covers, and is labelled so.
\end{definition}

Every statement below is labelled in this way. Theorem~\ref{thm:seed} and the frames of the later sections are thread
statements. The next paragraphs are weave statements, and their ranges are stated: ``every thread'' always means every
thread in the stated range: every word where a proof is stated, and a stated bound on the word length where a census is.

\paragraph{Three, from parity.}
Reduce the records modulo $2$. The group $\SL(2,\mathbb{F}_2)$ acts on the three non-zero vectors of $\mathbb{F}_2^2$, the
three \emph{parities} $(1,0),(0,1),(1,1)$, as the symmetric group $S_3$. Nothing is put in to make three: $2^2-1=3$.

\begin{lemma}[the parity lemma]\label{lem:parity}
With the grammar's action $(a,b)\mapsto(a,a+b)$, $L$ fixes the parity $(0,1)$ and exchanges the other two; $R$
fixes $(1,0)$; the swap $P:(a,b)\mapsto(b,a)$ of the two records (not a move of
the ruled grammar; see below) fixes $(1,1)$; the sign $-I$ fixes all three. Any two of $L$, $R$ and the swap generate $S_3$ on the three
parities. For $\phi\in\SL(2,\Z)$, $\tr\phi$ is odd if and only if $\phi$ acts on the parities as a $3$-cycle.
\end{lemma}
\begin{proof}
Modulo $2$, $L$ and $R$ are the transvections $\begin{psmallmatrix}1&1\\0&1\end{psmallmatrix}$ and $\begin{psmallmatrix}1&0\\1&1\end{psmallmatrix}$, which are transpositions of the three
non-zero vectors; two distinct transpositions generate $S_3$. An element of $\SL(2,\mathbb{F}_2)$ of order three has
characteristic polynomial $x^2+x+1$, hence trace $1$; the identity and the transpositions have trace $0$.
\end{proof}

The lemma is exact for every word. The record also checked it on all $988$ signed words to length eight ($448$ of odd
trace; B1600). Thus the three parities are cycled by every thread of odd trace, with no thread chosen, and a thread of
even trace fixes at least one of them.

\paragraph{The common point.}
For a thread, take the traces $(x,y,z)=(\tr a,\tr b,\tr ab)$ of the fibre's holonomy at the geometric point. The
Fricke maps of $L$ and $R$ act on the parabolic surface $x^2+y^2+z^2=xyz$ (commutator trace $-2$, the puncture), and on this
surface their only common fixed point is the origin $(0,0,0)$, the \emph{quaternion point}: there the holonomy sends
$a\mapsto i$ and $b\mapsto j$ in the quaternion group $Q_8\subset\mathrm{SU}(2)$ (B1600, W2). We call it the
\emph{common point}. It is the point at which every odd-trace thread meets the same finite picture.

The common point is also the geometry modulo three. For a thread with odd trace the ideal generated by $(x,y,z)$ in the
ring of integers of $\mathbb{Q}(x,y,z)$ is, on every thread to word length eight, one prime of residue degree one above
$3$, of norm $3$ and multiplicity one. The holonomy reduces modulo that prime to the quaternion point, and the cover
that forces the binary tetrahedral group $2T=\SL(2,\mathbb{F}_3)$ is the thread's principal congruence cover at it. For
even trace the ideal is a power of a prime above $2$, and the common point is the trivial representation (B1601).

\begin{censuslaw}[the geometry modulo three, B1601]\label{law:modthree}
Call a \emph{geometry} the bundle of a thread up to the symmetries that preserve it (cyclic rotation and the exchange of $L$
and $R$), so that each is computed once. On all $37$ geometries of word length at most eight ($16$ of odd trace, $21$ of even trace), an odd-trace thread has the
prime above $3$ described above, and an even-trace thread has no prime above $3$. The linearized lemma behind it is
proved: the Fricke maps' linear parts at the common point are quarter turns about the $x$- and $y$-axes of the three
parity lines, generating the rotations of the cube, and a word has odd trace if and only if its linear part is a
third turn about a body diagonal, an axis isotropic for $x^2+y^2+z^2$ exactly in characteristic three. That odd trace
\emph{always} gives this prime, and even trace never, is not proved for all words.
\end{censuslaw}

\paragraph{The parity triplet.}
The moves act on the three parities, and the common point carries $2T$. Together they give a triplet.

\begin{theorem}[the parity triplet, T2]\label{thm:triplet}
\begin{enumerate}
\item \textbf{(Parity.)} The moves act on the three non-zero parities of the two records as $S_3$, and a thread has odd
  trace if and only if it acts as a $3$-cycle (Lemma~\ref{lem:parity}; every word).
\item \textbf{(The common point.)} For every odd-trace word, there is an element $\tau\in2T\setminus Q_8$ conjugating
  $(i,j)$ to the images of the thread's generators, and the group generated by $i,j,\tau$ is $2T$, of order $24$. The group
  acts on the three parity lines as the standard three-dimensional irreducible representation of $A_4=2T/\{\pm1\}$
  (character $(3,-1,0,0)$, norm one). For an even-trace word the conjugator lies in $Q_8$ if the word acts trivially on the
  parities, and in the binary octahedral group $2O$ outside $2T$ if it acts as a transposition; the parity lines split.
  The arithmetic of the common point, a prime above $3$, is Census law~\ref{law:modthree}.
\item \textbf{(The zero modes.)} On the weave's surface $M_{1,2}$ (taken as a reading, Table~\ref{tab:inputs}), let $V$ be
  the parity-twisted spin space. The group generated by the lifts of $L$ and $R$ on $V=T\oplus\bar T$ has order $96$. The
  Hodge--Riemann form on $V$ is invariant under $L$, $R$ and the sign, has signature $(3,3)$ and is positive definite on
  $T$. The punctured fibre with this local system has three holomorphic zero modes, one per parity, and they span $T$.
\end{enumerate}
\end{theorem}
\begin{proof}[Proof (records B1600, B1606, B1624)]
(1) is Lemma~\ref{lem:parity}. (2) The Fricke maps fix the common point, so for a word $\varphi$ the representation $\rho\circ\varphi$ of the free
group, with $\rho(a)=i$ and $\rho(b)=j$, has the character of $\rho$. It is irreducible, so $\rho\circ\varphi=g\rho g^{-1}$
for some $g\in\mathrm{SU}(2)$, determined up to sign. Then $g$ normalises $Q_8$, so $g$ lies in $2O$, the normaliser of
$Q_8$ in $\mathrm{SU}(2)$. The class of $g$ in $2O/Q_8\cong S_3\cong\mathrm{Out}(Q_8)$ is the action of $\varphi$ on the parities.
So $g\in2T\setminus Q_8$ exactly when the trace is odd (Lemma~\ref{lem:parity}), $g\in Q_8$ when $\varphi$ is trivial
modulo $2$, and $g\in2O\setminus2T$ for a transposition; for odd trace $\langle i,j,g\rangle=2T$. The record checks this with
exact rational quaternion matrices for every word of length at most six (B1600); the first attempt with irrational entries
failed an equality test because of an error in the instrument, caught by an impossible group order and corrected before
anything was read off. The triplet is then a statement of representation theory:
$h^1(F;\chi_\varepsilon)=1$ for each non-zero parity $\varepsilon$ (the twisted cohomology of the punctured torus), and
$A_4=(\mathbb{Z}/2)^2\rtimes\mathbb{Z}/3$ permutes the three non-trivial characters of $(\mathbb{Z}/2)^2$ cyclically,
which is the standard irreducible three. For an involution the representation splits. (3) Here $V=\bigoplus_p H^1(F;\chi_p\otimes\rho_Q)$, with $\rho_Q(a)=i$, $\rho_Q(b)=j$ the
common point's quaternion representation and $\chi_p$ the parity character. A move $m$ lifts through the element $g$ of
$2O$ with $\rho_Q\circ m=g\rho_Q g^{-1}$: $g_L=(1+i)/\sqrt2$ for $L\colon a\mapsto a,\ b\mapsto ab$, and
$g_R=(1-j)/\sqrt2$ for $R\colon a\mapsto ab,\ b\mapsto b$. The cup product on the
$\mathbb{Z}/2$-orbifold presentation of the fibre group, by an independent implementation, gives the form and its
signature; the zero modes are $F=(f(z),f(z+\tau))\,dz$ with $f=\theta_3(z;2\tau)/\sqrt{\theta_1(z;\tau)}$ and the three
parity twists $F_p=C_pF$, with $C_p$ the unit of $Q_8$ attached to the parity $p$ and the square root continued along
the segments that represent $a$ and $b$, so that $F$ has the monodromy of $\rho_Q$ (W21).
\end{proof}

\noindent\emph{Scope.} Parts (1) and (2) are theorems: weave statements for every word, with no bound on the length. The
prime above $3$ of Census law~\ref{law:modthree} is a census. Part (3) holds on the
surface $M_{1,2}$, whose identification with the weave's object is a reading made within the program, of which only the
Euler characteristic is verified ($\chi=1/12$ by two routes; B1624). The identification of the zero modes with
left-handed matter is a further reading and is not used here. The index of the puncture (odd, never zero, and $\pm3$
once the parity grading is kept, B1606) belongs to Section~\ref{sec:three}. The triplet is a flavour structure: the
record grades the flavour three as derived on the weave and the gauge three as a selection, and the paper claims the first
only. The verification of (1)--(3) in the record is B1600 and B1606.

\paragraph{What is classical.}
A triplet of $A_4$ on three lines permuted by an order-three element is the standard flavour triplet of discrete flavour
models; we cite the use and not the construction~\cite{MaRajasekaran2001,AltarelliFeruglio2005}. What the weave adds is that the three lines, the group and the
irreducibility come from the records' parities and the common point, with no choice of group: proved for every odd-trace
word.

\begin{remark}[the moves that are not there]
The swap of the two records and the inner automorphisms of the free group by $a$ and $b$ are not moves of the
grammar adopted here. This matters in Section~\ref{sec:modular}: the group $\langle L,R\rangle\cong B_3$ meets
$\mathrm{Inn}(F_2)$ only in $\langle\Delta^4\rangle$, so the weave's group is the braid group, not its extension by the
point-pushes. Results of the program that take the inner automorphisms as symmetries hold on branches that add them, and
say so (B1634, B1640).
\end{remark}

\paragraph{One example.}
The root's own triplet is the running example of the paper. Its trace $3$ is odd, so $LR$ cycles the three parities, its
common point is the quaternion point, and its three parity lines form the standard triplet of $A_4$. It is one of the
odd-trace threads of Theorem~\ref{thm:triplet}, with no property, for the triplet, that the others lack.
